\section{Introduction}

A classical intuition in auction theory, formalized by \citet{bulow1996}, is that bringing additional bidders to the table matters more than any sophisticated bidding rule. A set-aside that restricts competition to small and medium-sized enterprises (SMEs) is in direct tension with that intuition: it bars some bidders entirely, and it does so by identity. Yet when a March~2018 S\~ao Paulo reform extended the state's SME-only rule to medical supplies, SME participation in BEC auctions almost doubled---from 0.55 to 1.22 firms per auction in pharmaceuticals, from 0.94 to 1.87 in non-pharmaceuticals---and winning prices rose 10--11 percent. More bidders did not soften the price cost. They arrived, and the price cost arrived with them.

This paper is about the half of that fact that the reduced-form literature cannot see. A \citet{gelbach2016} decomposition of the 10--11 percent price rise against observable entry and winner-composition channels assigns those channels five percent of the effect. The remaining 95 percent does not flow through who enters, who wins, or how concentrated the bidder pool gets. It flows through how each bidder bids. Bid distributions shift toward the reference price when the rule binds, and they shift in ways that firm counts, winner identities, and geographic sorting cannot account for. The \emph{within-auction} response of surviving bidders to the exclusion of their rivals is the policy's dominant channel. I call it \emph{sheltered bidding}, and quantify it structurally.

\paragraph{Setting and identification.} SME set-asides are the most widely deployed procurement-preference instrument in the world. The United States Small Business Administration has operated them since the 1950s, the European Union catalogues them in more than one hundred jurisdictions, and Brazilian Law~14{,}133/2021 made them a procurement default \citep{bastos2018, thai2017, bosio2022}. Existing empirical evidence places the price cost of such rules at 10--15 percent \citep{marion2007, krasnokutskaya2011, athey2013, nakabayashi2013}; what that evidence has not been able to settle is \emph{why}. When a rule restricts who can bid, the set of admissible firms moves and the competitive pressure each surviving firm faces moves; price data alone cannot separate the two margins.

On February~6, 2018, a S\~ao Paulo public buyer procured fifty thousand units of furosemide 40\,mg---a generic antihypertensive, item 110639 in the BEC catalog---through an open Preg\~ao. Ten firms bid, and the clearing price settled twelve percent below the buyer's reference. Eight months later, after the state attorney's office reversed an administrative exemption that had kept medical supplies outside the SME-only rule since 2014, a comparable purchase by the same buyer attracted three small firms, and the winning unit price cleared just above the reference. The price gap is about fourteen percent. It is neither an outlier nor a coincidence: it is the average direction of more than a hundred thousand item-transactions in a matched eighteen-month window. The March~2018 reversal was a legalistic reinterpretation internal to PGE-SP, triggered by an isonomy argument and untied to price levels, competition patterns, or supplier complaints specific to medical supplies. The timing was bureaucratic, and the switch was plausibly exogenous to procurement outcomes. This is the natural experiment the paper exploits.

A difference-in-differences against the 76 product groups that were already under the SME-only rule delivers the reduced-form headline. Winning prices rose 10--11 percent in nominal terms (7--10 percent IPCA-deflated), valid bid counts fell 7--8 percent, participating firm counts fell 10--11 percent, and winner-to-buyer distance rose six to fourteen kilometers. The estimates survive staggered-DiD corrections \citep{callaway2021, sun2021, goodmanbacon2021, borusyak2024}, a formal spillover bound on non-SME displacement (below one percent of the headline), matched-employer-employee validation against RAIS, wild cluster bootstrap, synthetic control, and four placebo cutoffs. The Gelbach decomposition isolates the observable channels; they explain five percent of the effect. The residual 95 percent---rather than the headline 10--11 percent---is the object the structural analysis must recover.

\paragraph{Structure.} The structural layer specifies the BEC Preg\~ao as an asymmetric independent-private-values auction of the descending-clock English-reverse format. Under \citet{haile2003}, drop-out prices of losing bidders point-identify their private costs, so the cost distribution is recovered non-parametrically without inverting a bid function or imposing parametric structure on bidding. An auction-level unobserved heterogeneity component is shrunk out along \citet{krasnokutskaya2011} lines; entry is treated in the \citet{atheyseira2011} spirit, with observed post-policy counts serving as the equilibrium response to the restricted regime.

Three findings anchor the structural decomposition. First, the intensive margin of bidding accounts for 75--85 percent of the simulated set-aside price effect, stable across three cost-distribution estimators (all-bidders, Turnbull left-censoring, losers-only) and across pharmaceutical and non-pharmaceutical classes. A Bayes-Nash equilibrium simulation translates the 10--11 percent reduced-form rise into $+0.24$ of the reference price in non-pharma and $+0.34$ in pharma, with bootstrap 95 percent intervals of $[0.19, 0.29]$ and $[0.25, 0.36]$. Second, endogenous SME entry absorbs 60--70 percent of what the same exclusion would deliver at a frozen SME pool: holding the pool at its pre-policy size raises the simulated price effect to 162 percent of observed in non-pharma and 170 percent in pharma. The induced entry response is the government's partial insurance, not the source of the markup. Third, welfare decomposes into an allocative wedge and a tax-distortion term through the marginal cost of public funds \citep{ballard1985}; the total welfare cost of the set-aside is 29 percent of procurement value in non-pharma (95 percent CI 21--35) and 45 percent in pharma (95 percent CI 35--56).

\paragraph{Policy.} A counterfactual 10 percent SME price preference---of the type the U.S. SBA has operated for decades and that Brazilian procurement law already contemplates---delivers comparable SME win-rate gains at essentially zero price cost in thick, standardized markets. Non-SMEs remain admissible and keep competing; SMEs receive protection only to the extent that their bids fall within 10 percent of the non-SME minimum. The implicit social welfare weight required to prefer the full set-aside over the preference rule is 2.4 in non-pharma and 3.0 in pharma, above the 1.2--1.5 range typical of Brazilian cash-transfer programs (Bolsa Fam\'{\i}lia, Aux\'{\i}lio Brasil) and far above the utilitarian benchmark. In thin, heterogeneous markets the ranking can reverse---a boundary result that identifies where bidder-exclusion policies cease to deliver proportional returns, not a fragility of the preference-rule conclusion where it dominates.

\paragraph{Contribution.} The paper sits in direct dialogue with the structural procurement-preference literature of \citet{marion2007}, \citet{krasnokutskaya2011}, and \citet{atheyseira2011}, where the setting is \emph{partial} preference---SMEs and non-SMEs both remain active in the auction with differential treatment. In a partial-preference setting the intensive and extensive margins are entangled by construction: the policy changes both what is profitable to bid and who finds it profitable to enter, and each type's response to each margin is embedded in the other's residual. The Brazilian reform is \emph{total} exclusion: non-SMEs leave the auction entirely, the SME pool re-equilibrates under a cleanly smaller bidder set, and the two margins can be separated without imposing parametric structure on the response. Four contributions follow.

\emph{First}, I isolate the within-auction channel structurally under total bidder exclusion, where drop-out data alone---rather than a bid-function inversion---identifies the cost distribution. The partial-preference settings of \citet{krasnokutskaya2011} and \citet{atheyseira2011} cannot decouple the two margins on the same footing.

\emph{Second}, I deliver the first developing-country evidence on an SME set-aside paired with both a causal reduced-form estimate and a structural margin decomposition, using a legally-triggered quasi-experiment and bidder-level drop-out data the prior literature---largely confined to U.S. federal and timber settings---did not have access to.

\emph{Third}, I show that welfare accounting for bidder-exclusion rules must incorporate endogenous entry as a self-insurance mechanism; a fixed-pool calculation, the standard reduced-form shortcut, inflates the estimated welfare loss by approximately two-thirds.

\emph{Fourth}, I identify and quantify a specific policy instrument---a 10 percent price preference---that Pareto-dominates the current Brazilian set-aside in thick, standardized markets, and map the market conditions under which this dominance does and does not hold.

The paper proceeds as follows. Section~\ref{sec:institutional} describes the BEC setting and the legal reversal. Section~\ref{sec:data} describes the data. Section~\ref{sec:rf} presents the reduced-form evidence and the Gelbach residual. Section~\ref{sec:bridge} states the questions the reduced form leaves open. Sections~\ref{sec:model}--\ref{sec:identification} specify the asymmetric IPV model and the identification strategy. Section~\ref{sec:results} reports the structural decomposition and the partial-insurance finding. Section~\ref{sec:welfare} computes welfare. Section~\ref{sec:pareto} compares the set-aside against a 10 percent price preference. Section~\ref{sec:discussion} discusses the goods-characteristic lesson. Section~\ref{sec:conclusion} concludes.
